\chapter{Cierre: de la teoría al código}
\label{cap:cierre}

\section{Un mapa de todo lo recorrido}

Llegamos al final con una convicción que vale la pena hacer explícita: \textbf{los nueve bloques de este curso no son nueve temas sueltos, sino una sola historia}. Empezamos describiendo la actividad económica con funciones; las pusimos a interactuar para hallar equilibrios; medimos el cambio con derivadas y lo acumulamos con integrales; usamos ese cálculo para optimizar; llevamos la optimización al terreno estratégico del oligopolio; y proyectamos todo hacia el futuro con la evaluación de inversiones. Cada pieza se apoya en la anterior.

La siguiente tabla resume el puente entre cada concepto económico y la herramienta computacional que lo hace operativo en el laboratorio:

\begin{center}
\renewcommand{\arraystretch}{1.25}
\begin{tabular}{p{5.6cm} p{3.4cm} p{4.2cm}}
\toprule
\textbf{Concepto económico} & \textbf{Capítulo} & \textbf{Herramienta} \\
\midrule
Datos, variables y visualización & \ref{cap:datos} & \texttt{pandas}, \texttt{matplotlib} \\
Funciones (demanda, costo, ingreso) & \ref{cap:funciones} & \texttt{sympy}, \texttt{matplotlib} \\
Equilibrio de mercado & \ref{cap:equilibrio} & \texttt{sympy.solve} \\
Insumo-producto (Leontief) & \ref{cap:leontief} & \texttt{numpy.linalg} \\
Análisis marginal & \ref{cap:marginal} & \texttt{sympy.diff} \\
Elasticidad & \ref{cap:elasticidad} & \texttt{sympy}, \texttt{pandas} \\
Excedentes & \ref{cap:integrales} & \texttt{sympy.integrate} \\
Optimización & \ref{cap:optimizacion} & \texttt{sympy}, \texttt{PuLP} \\
Duopolio de Cournot & \ref{cap:duopolio} & \texttt{sympy.solve} (sistema) \\
Evaluación de inversiones & \ref{cap:inversiones} & \texttt{numpy\_financial} \\
\bottomrule
\end{tabular}
\end{center}

\section{Qué aporta cada mitad}

Conviene insistir en la división del trabajo entre el economista y la máquina, porque es el sentido último de la materia:

\begin{itemize}
  \item \textbf{La computadora} hace el cálculo: deriva, integra, invierte matrices, resuelve sistemas, descuenta flujos. Es rápida, exacta y no se cansa.
  \item \textbf{El economista} hace lo que la máquina no puede: \emph{plantear} el problema correcto, \emph{elegir} qué función representa la realidad, \emph{interpretar} si un resultado tiene sentido y \emph{decidir} en consecuencia.
\end{itemize}

Un VAN positivo no decide por vos; te informa. Una elasticidad de $-0.7$ no fija el precio; te dice que tenés margen para subirlo. El juicio económico —saber qué preguntar y cómo leer la respuesta— es insustituible, y es exactamente lo que esta materia busca entrenar.

\section{Para seguir}

Este material es una puerta de entrada, no una enciclopedia. Quien quiera profundizar encontrará en la microeconomía de \cite{varian}, \cite{pindyck} y \cite{nicholson}, en la matemática económica de \cite{chiang} y \cite{sydsaeter}, en las finanzas de \cite{brealey} y \cite{sapag}, y en la macro de \cite{blanchard}, desarrollos mucho más extensos de cada tema. Pero la lógica de fondo —traducir un problema de gestión a una estructura económica, resolverla con la herramienta adecuada e interpretar el resultado— es la que ya tenés, y la que vas a usar una y otra vez.

\begin{ideabox}
\textbf{La síntesis del curso, en una línea.} La teoría económica te dice \emph{qué} mirar; los métodos cuantitativos te dicen \emph{cuánto} vale; y el criterio profesional te dice \emph{qué hacer} con esa respuesta.
\end{ideabox}
